\begin{thebibliography}{9}
\addcontentsline{toc}{chapter}{Sources and references}
\bibitem{nist45} Scott Rose, Oliver Borchert, Stu Mitchell and Sean Connelly.
\emph{Zero Trust Architecture}. NIST SP 800-207, 2020. Official publication
overview inspected for resource authentication/authorization and network-location
trust boundary.
\url{https://csrc.nist.gov/pubs/sp/800/207/final}.
\bibitem{injection45} OWASP Gen AI Security Project.
\emph{LLM01:2025 Prompt Injection}. Types and mitigation sections inspected.
\url{https://genai.owasp.org/llmrisk/llm01-prompt-injection/}.
\bibitem{boundary45} OWASP Gen AI Security Project.
\emph{LLM07:2025 System Prompt Leakage}. Mitigation/privilege-boundary guidance
inspected. \url{https://genai.owasp.org/llmrisk/llm072025-system-prompt-leakage/}.
\end{thebibliography}

\section*{Evidence scope}
Access checked 17 October 2026. NIST's publication overview was inspected for
the no-implicit-network-trust principle; this is not a full implementation
audit against every SP 800-207 requirement. OWASP's current LLM01:2025 and
LLM07:2025 pages were inspected for mitigation and external authorization
boundaries. Their recommendations are source leads, not authorization to test
live systems. No penetration test was performed. The capability registry,
namespace proof, quota oracle, local adversaries and mechanism figures are
original teaching constructions. They assume trusted serialized
host code and do not establish cryptographic or hardware isolation.
